\section{Conclusion}
\label{sec:conclusion}

\paragraph{Conclusion.}
Credit that rewards agents for being relied on rewards attention, not accuracy, and learning agents respond by exaggerating. We proposed truthful verified credit (\tvc{}), which gives credit only for verified value, adds a bonus for every checked claim, and fines audited false claims. The fine needed to outweigh a lie's outside gain grows as $1/p$ in the audit rate $p$. Across both experiments, \tvc{} nearly eliminates exaggeration and keeps agents reporting. Its welfare is comparable to popularity credit when receivers can decode exaggeration cheaply, and well above both baselines when receivers act on claims. Lying stays low at every prosocial weight and collapses close to the theoretical fine threshold. \tvc{} is called for when credit follows reliance, lying is profitable outside the credit, lying harms others, and a sample of claims can be audited. When agents weigh the others' payoffs almost as much as their own, a shared return alone comes close.

\paragraph{Limitations.}
\tvc{} elicits truth, not usefulness. In the bandit its welfare matches popularity credit but does not exceed it, because calibrated receivers decode exaggeration cheaply, and designs that raise exploration reach higher welfare without making reports honest. The fine rule is derived for a sender that maximises expected reward in a single round. It also assumes that checks reveal the truth exactly, that audits are independent, and that receivers never follow a claim that cannot help them. Our testbeds are simulated, and the outside gain of a lie is calibrated from short runs, not an emergent property. The bonus accounts for most of the credit given in the congestion setting, yet removing it changes little there. Future work includes noisy verification, the cost of audits, and fleets with real data.
